% !PRINTED 19
\weberpage{19}

If then $\Phi(x,y)$ is an integral or fractional rational function
of $x$ and $y$, and $y$ depends on $x$ through equation (1), then
\weq{2}{\Phi(x,y)\,dx, \qquad \int \Phi(x,y)\,dx}
are the \gesperrt{algebraic differentials and integrals} belonging to the
curve $F$.

For example, if $f(x)$ is a function of the third or fourth degree
in $x$, then
\[ y^2 - f(x) = 0 \]
is the equation of a curve of the third or fourth degree, to which
belong the \gesperrt{elliptic differentials and integrals} affected with the
irrationality $\sqrt{f(x)}$.

We shall here call all curves whose differentials and integrals
are reducible to elliptic ones \gesperrt{elliptic curves}. As the
example shows, curves of the third and fourth degree occur among
them.

Conics do not belong to the elliptic curves, because, when an
equation of the second degree subsists between $x$ and $y$, $x$
and $y$ can be represented as rational functions of a parameter
$t$, whereby the algebraic differential can be reduced to a
rational one in $t$. Such curves are called \gesperrt{rational
curves}.

We shall see that all curves of the third degree without a double
point or point of regression belong to the elliptic ones. Curves
of higher degree can belong to them only when they have a certain
number of singular points. This is a fundamental chapter in the
general theory of algebraic functions, which does not belong to
the plan of this work\footnote{The most important papers
concerning this subject are:

\gesperrt{Riemann}, Theorie der \gesperrt{Abelschen} Funktionen [Crelles
Journal, Bd.\ 54 (1857)]. Mathematische Werke, 2.\ Aufl., S.\ 88,
487. Nachträge, edited by \gesperrt{Noether} and
\gesperrt{Wirtinger} (1902).

\gesperrt{Aronhold}, Monatsberichte der Berliner Akademie of
25 April 1861.

\gesperrt{Clebsch}, Über die Anwendung der \gesperrt{Abelschen} Funktionen
in der Geometrie (Crelles Journal, Bd.\ 63).

\gesperrt{Clebsch}, Über diejenigen ebenen Kurven, deren
Koordinaten rationale Funktionen eines Parameters sind (ibid.,
Bd.\ 64).

\gesperrt{Clebsch}, Über diejenigen Kurven, deren Koordinaten
sich als elliptische Funktionen eines Parameters darstellen
lassen (ibid., Bd.\ 64).}.
